\homework{10}{Thu 06 Apr 2023 23:55}{}

§21 2,4, §22 1,2, §16 1,3

\begin{itemize}
	\item [21.2] Let $x_1, \ldots, x_k$ be vectors in $\mathbb{R}^n$. Show that
		\[
			V(x_1, \ldots, \lambda x_i, \ldots, x_k) = |\lambda| V(x_1, \ldots, x_k)
		.\] 
\begin{proof}
	Write $X = [x_1, \ldots, x_k]$ to be an $n \times  k$ matrix. Then
	\[
		X E^T = [x_1, \ldots, \lambda x_i, \ldots, x_k]
	.\] 
	where $E$ is the $k \times  k$ matrix with the $i$-th entry replaced to $\lambda$. Obviously $\text{det}E = |\lambda|$.

	Now
	\[
		F(X E^T) = \text{det}(E X^T X E^T) = (\text{det} E)^2 F(X)
	.\] 
	Hence $V(XE^T) = \text{det}E \,V(X) = |\lambda| V\left(X \right) $.
\end{proof}

\item [21.4] Use Pythagorean Theorem to verify 
	\[
		\|a\|^2 \|b\|^2 - \left\langle a,b \right\rangle ^2 = \|a \times b\|^2
	.\] 
\begin{proof}
	\begin{align*}
		\|a \times b\|^2
		&=
		\text{det}
		\begin{bmatrix}
			a_2 & b_2 \\
			a_3 & b_3 \\
		\end{bmatrix}
		^2
		+
		\text{det}
		\begin{bmatrix}
			a_1 & b_1 \\
		a_3 & b_3 \\
		\end{bmatrix}
		^2
		+
		\text{det}
		\begin{bmatrix}
			a_1 & b_1 \\
			a_2 & b_2 \\
		\end{bmatrix}
		^2\\
		&= \sum_{[I]} \text{det}^2 X_I \\
		&= V(X)\\
		&= \|a\|^2\|b\|^2 - \|a\|^2 \|b\|^2 \cos^2 \theta\\
		&= \|a\|^2 \|b\|^2 - \left\langle a,b \right\rangle ^2 
	.\end{align*}

\end{proof}

\item [22.1] Let $A$ be open in $\mathbb{R}^k$, let $\alpha: A \to \mathbb{R}^n$ be of class $C^r$, let $Y = \alpha(A)$. Suppose $h: \mathbb{R}^n \to \mathbb{R}^n$ is an isometry. Let $Z = h(Y)$ and let $\beta = h \circ \alpha$. Show that $Y_\alpha$ and $Z_\beta$ have the same volume.
\begin{proof}
	By definition, $V(Y_\alpha) = \int_A V(D\alpha)$, and $V(Z_\beta) = \int _A V(D \beta)$. Then it suffices to show $V(D \alpha) = V(D \beta)$.
	 \begin{align*}
		F(D \beta)
		&=  F(D(h \circ \alpha))\\
		&= F(Dh D\alpha) \\
		&= \text{det}(Dh) F(D \alpha)\\
		&= F(D \alpha)
	.\end{align*}
	Hence $V(D \alpha) = V(D \beta)$.
\end{proof}
\item [22.2] Let $A$ be open in $\mathbb{R}^k$, let $f: A \to \mathbb{R}$ be of class $C^r$. Let $Y$ be the graph of $f$ in $\mathbb{R}^{k+1}$, parameterized by the function $\alpha: A \to \mathbb{R}^{k+1}$ given by $\alpha(x) = (x, f(x))$. Express $v(Y_\alpha)$ as an integral.
\begin{proof}
	Write $h$ to be the identity on $\mathbb{R}^k$, then we have $\alpha(x) = (h(x), f(x))$. Then
	
	\[
		D \alpha (x) =
		\begin{bmatrix}
			I \\
			Df(x) \\
		\end{bmatrix}
	\]
	.
	Now $D\alpha(x)$ is a $(k+1)$ by $k$ matrix. Consider $D\alpha(x)|_{I_i}$, the removal of $i$-th row from $D\alpha(x)$. What is its determinant squared? If $i =k$, then $D\alpha(x)|_{I_k} = I$, the $k \times k$ identity, so the determinant squared is $1$.

	If $i < k$, then note that the $i$-th colomn have zero entries on all but the last row. To compute the determinant, remove the $i$-th column and last row and the remaining is a $(k-1) \times (k-1)$ identity matrix. Then we have
	\[
		\text{det} D\alpha(x)_{I_i} = (-1)^{k + i} \frac{\partial f}{\partial x_i} \text{det}I
	.\] 
	Hence
	\[
		\text{det}^2 D\alpha(x)_{I_i} = \left( \frac{\partial f}{\partial x_i} \right) ^2
	.\] 
		
	Now by the Pythagorean Theorem,
	\[
		V(D \alpha(x)) = \left( \sum_{i \le k} \text{det}^2 D\alpha(x)_{I_i} \right) ^{\frac{1}{2}} = \left(\left( \frac{\partial f}{\partial x_1} \right)^2  + \left( \frac{\partial f}{\partial x_2}  \right) ^2 + \ldots + \left( \frac{\partial f}{\partial x_{k}}  \right)^2 + 1 \right) ^{\frac{1}{2}}
	.\] 

	Put everything together,
	\begin{align*}
		v(Y_\alpha)
		&= \int _A V(D \alpha) \\
		&= \int _A \sqrt{1+\sum_{i=1}^{k} \left( \frac{\partial f}{\partial x_i}  \right) ^2} 
	.\end{align*}

	

\end{proof}

\item [16.1] Prove that the function $f$ of lemma 16.1 is of class $C^{\infty } $ as follows: Given any integer $n\ge 0$, define $f_n: \mathbb{R} \to \mathbb{R}$ by the equation
	\[
		f_n(x) = \begin{cases}
			\frac{e^{-\frac{1}{x}}}{x^n} & x>0\\
			0 & x \le 0
		\end{cases}
	.\] 
\begin{proof}
	\begin{itemize}
		\item First show continuity. Set $y = \frac{1}{x}$ and assume $x>0$.
			\begin{align*}
				f_n(y) &= e^{- y} y^{n} \\
				&= \frac{y^n}{1+ y + \ldots + y^n / n! + \ldots} \\
				&\le \frac{y^n}{y^{n+1}/(n+1)!} \\
				&= \frac{(n+1)!}{y} \to 0 \quad \text{ as } y \to \infty 
			.\end{align*}
			Hence $f_n(x) \to 0$ as $x \to 0^+$. Obviously $f_n$ continuous from the left at $0$, so $f_n$ continuous at $0$.
		\item Next show differentiability. It suffices to show that $f_n$ has a right derivative at $0$ and the derivative is $0$. Let $h>0$. We have as $h \to 0$,
			\[
				\frac{f_n(h) - f_n(0)}{h} = e^{- \frac{1}{h}} h^{- n - 1} = f_{n+1}(h) \to 0
			.\] 
		\item We establish identity in $(c)$, For $x \le  0$ the claim is trivial. Now we assume $h > 0$. Then we can just do calculus.
			\[
			\frac{\partial f_n(x)}{\partial x}  = \frac{1}{x^2} f_n(x) - n x^{-n-1} e^{-\frac{1}{x}} = f_{n+2}(x) - n f_{n+1}\left( x \right) 
			.\] 
		\item Finally show $f$ is $C^{\infty }$. Suppose $f_n$ is $r$-times differentiable and $r \in \mathbb{N}$. Then by the result above, either $f_{n+1} $ or $f_{n+2}$ is $r-1$-times differentiable. Then we trace the derivative of both of them, and repeat the process. In each step, we know one of $f_{n + k} , \ldots, f_{n+ 2 k} $ must be $r - k$-times differentiable. This process must terminate, and we know one of $f_{n+r}, \ldots, f_{n+ 2r}$ is nondifferentiable. But this is a contradiction since we have proven that $f_m$ is differentiable for all $m$.
	\end{itemize}
\end{proof}
\item 3. (a) Let $S$ be an arbitrary subset of $\mathbf{R}^n$; let $x_0 \in S$. We say that the function $f: S \rightarrow \mathbf{R}$ is differentiable at $x_0$, of class $C^r$, provided there is a $C^r$ function $g: U \rightarrow \mathbf{R}$ defined in a neighborhood $U$ of $x_0$ in $\mathbf{R}^n$, such that $g$ agrees with $f$ on the set $U \cap S$. In this case, show that if $\phi: \mathbf{R}^n \rightarrow \mathbf{R}$ is a $C^r$ function whose support lies in $U$, then the function
\[
h(\mathbf{x})= \begin{cases}\phi(\mathbf{x}) g(\mathbf{x}) & \text { for } \mathbf{x} \in U, \\ 0 & \text { for } \mathbf{x} \notin \text { Support } \phi,\end{cases}
\]
is well-defined and of class $C^r$ on $R^n$.
\begin{proof}
	Since $\varphi$ and $g$ are both $C^r$, their product is $C^r$. Since support of $\varphi$ lies in $U$, $\varphi = 0$ outside $U$, so $\varphi g = 0$ outside $U$. In fact, the second line in the definition of $h(x)$ is redundant.
\end{proof}

(b) Prove the following:
Theorem. If $f: S \rightarrow \mathbf{R}$ and $f$ is differentiable of class $C^r$ at each point $x_0$ of $S$, then $f$ may be extended to a $C^r$ function $h: A \rightarrow R$ that is defined on an open set $A$ of $\mathbf{R}^n$ containing $S$.
[Hint: Cover $S$ by appropriately chosen neighborhoods, let $A$ be their union, and take a $C^{\infty}$ partition of unity on $A$ dominated by this collection of neighborhoods.]
\begin{proof}
	At each $x \in S$ there exists some neighborhood $U_x$ of $x$, and function $g_{U_x}: U_{x} \to \mathbb{R}$ as in the definition of differentiability of $f$. Then the collection $\mathcal{C} = \{U_x: x \in S\} $ is an open collection that covers $S$. Then there exists a partition of unity $\varphi_n$ of class $C^r$ and covers $S$, with compact support and dominated by $\{U_x\} : x \in S$. For each $\varphi_n$ there is some associated $U_n$ and $g_n: U_n \to  \mathbb{R}$ of class $C^r$ where $g_n$ agrees with $f$ on $U \cap S$. Define
\begin{align*}
	h: A &\longrightarrow \mathbb{R} \\
	x &\longmapsto h(x) = \sum_{n=1}^{\infty} \varphi_n(x) g_n(x)
.\end{align*}
For any $x \in A$, there exists neighborhood $U$ and natural number $M$ such that $h(x) = \sum_{n=1}^M \varphi_n(x) g_n(x)$ for $x \in U$. Then we see that $h(x)$ is of class $C^r$ in the neighborhood $U$, since it is finite sum of $C^r$ functions. This is true for any $x \in A$, so $h$ is of class $C^r$.

Finally, if $x \in S$, then
\[
	h(x) = \sum_{n = 1}^\infty  \varphi_n(x) f(x) = f(x) \sum_{n=1}^{\infty} \varphi_n(x) = f(x)
.\] 
Therefore $h$ is the desired extension.
\end{itemize}
